\section{Discussion and limits}
\label{sec:discussion}

Search, not the check, removed the wrong answers. The model alone gave a wrong value for 4.3 held-out questions per run; with eight retrieved excerpts it gave one wrong answer in three runs, and that answer quoted a real passage, so the check passed it. On these questions the check did not make the shown answers more accurate.

What the check bought is a property of every answer shown: each number in it appears in a passage quoted word for word from the page it cites. A reader who doubts a value opens one page, finds a highlighted passage and confirms the number there, rather than searching the datasheet. Without the check, arm B's answers carry the same citations, but nothing confirms that the quote is on the page or that the number is in the quote. On this set B's citations happened to hold up; the check is what makes that true of every answer rather than of those examined.

The check alone paid for this property with recall: it withheld 7.3 of B's 33.3 correct held-out answers per run. Its rejections were mostly the model's own habits rather than unsupported claims. A three-word quote such as \enquote{\enquote{110}: 10 ms} is the right evidence and simply too short for the rule; \enquote{bits 3..1} is register notation the model wrote in its own words. Because these rejections are systematic, sending the reason back works: one revision recovered 20 of the 22 withheld answers and brought arm D to within 0.6 correct answers and 0.7 false refusals of B. The two answers it could not recover came from one case the check gets wrong. A field name taken from the question, \texttt{t\_sb[2:0]}, contains digits, and the check counts them as numbers the quote must contain. Exempting text copied from the question would fix this case, but it would also let an answer repeat a wrong value that the question itself contained.

\subsection{What the check does not catch}

The check confirms that a number is in a quote and the quote is on a page. It does not confirm that the quote answers the question. An answer can quote the wrong row, the wrong condition or the wrong column of a table and pass, and the BMA456 sensor time question shows that a datasheet can state two values for one quantity. A revision can also pass by dropping context: in the development run, one revision dropped the $\pm 8\,g$ range from an answer about the accelerometer's reset range, and it was still scored correct. The revision prompt forbids dropping the value the question asks for, but nothing in code enforces it, and a revision could in principle swap the requested value for another number from the same quote. The scorer would count a swapped value as wrong, but it cannot see dropped context, and the assistant prevents neither.

\subsection{What the assistant cannot read}

The assistant reads only the text layer. Graphs, timing diagrams and package drawings are invisible to it, and formulas extract as garbled text. In tables the column a value sat in is lost, so a row of minimum, typical and maximum values cannot tell the model which is which; two of the withheld held-out answers said as much. The number extractor reads \enquote{10,000} as 10 and 0, so a value written with a thousands separator can be neither verified nor scored, and one held-out candidate was dropped for that reason. The check compares numbers as written, so an answer converted from \qty{1}{\milli\ampere} to \qty{1000}{\micro\ampere} is withheld, and the prompt asks the model not to convert. Decimal commas and scanned datasheets without a text layer are not handled.

\subsection{How far the result reaches}

The held-out result rests on 50 questions, 40 of them answerable, on six datasheets from one manufacturer, answered by one model. The questions are on the same datasheets as the development set, so they test the system on new questions, not on new documents. The held-out set was drafted and checked by Claude, and Claude also answers the questions; a person checking each question against the PDF, as was done for the development set, would make the set independent of the model it measures. Three runs show how much the counts move between runs, at most four questions for any arm and category, but they are too few for an interval. The scorer undercounts correct answers when the unit is written against the number, which leaves three correct answers per run unplaced in arms B to D, and it marks a correct answer wrong when the question records only one of two values that the datasheet gives.
